% =============================================================================
% Subsection 12b: Oracle (upper-bound) latent-ability subset selection.
% Nests under the MCQ "Experiments" section; splice where the red
% [ORACLE SUBSET RESULT PENDING] placeholder sits in 12_mcq_experiments.tex
% (it carries \label{subsec:mcq-oracle}).
% Bib: refs/oracle_subset.bib
%
% FIGURES REFERENCED (\includegraphics{figures/<basename>}).
% Integrator: copy each basename into ./figures/ from the ABSOLUTE source path.
%   figures/oracle_curve_ifeval.png
%     <- /Users/arhant/Documents/EDLM/olmo-eval-full/AdaptiveTesting/Research/01_MCQ_ATLAS/data/oracle_subset_selection/oracle_curve_ifeval.png
%   figures/latent_axis_ifeval.png
%     <- /Users/arhant/Documents/EDLM/olmo-eval-full/AdaptiveTesting/Research/01_MCQ_ATLAS/data/oracle_subset_selection/latent_axis_ifeval.png
%   figures/oracle_curve_math.png
%     <- /Users/arhant/Documents/EDLM/olmo-eval-full/AdaptiveTesting/Research/01_MCQ_ATLAS/data/oracle_subset_selection/oracle_curve_math.png
%   figures/latent_axis_math.png
%     <- /Users/arhant/Documents/EDLM/olmo-eval-full/AdaptiveTesting/Research/01_MCQ_ATLAS/data/oracle_subset_selection/latent_axis_math.png
%
% TABLE SOURCE DATA (absolute):
%   Table (min-N-to-hold) <- .../data/oracle_subset_selection/README.md, results.csv
%
% PREAMBLE REQUIREMENTS (main document): \usepackage{graphicx}, \usepackage{xcolor},
%   \usepackage{booktabs}.
%
% Convention: this is a fully-AI subsection. Prose is wrapped in \aiadd{...}
%   (renders blue) and the subsection closes with an \aibullets{...} skeleton a
%   human can rewrite from. [CITATION NEEDED] markers render red. Every number
%   traces to oracle_subset_selection/results.csv and README.md.
% =============================================================================

\providecommand{\aiadd}[1]{\textcolor{blue}{#1}}
\providecommand{\aibullets}[1]{{\color{blue}#1}}

% =============================================================================
\subsection{Oracle latent-ability subset selection}
\label{subsec:mcq-oracle}

\aiadd{The diversity experiment in Section~\ref{subsec:mcq-diversity} showed that
choosing calibration models well saves models over random sampling, but only by
about $1.4$ to $1.6\times$. That leaves a ceiling question: how much could the
best possible subset save if we allowed it to use information that is available
only after calibration? We answer it with an oracle. We calibrate the 3PL bank on
the whole capped pool~\cite{li2026adaptive,chalmers2012mirt}, read each model's
fitted latent ability (its EAP $\theta$)~\cite{lord1980irt,embretson2000irt} from
that bank, and then pick subsets that either spread across the fitted ability
range or directly maximize the held-out correlation. Both choices use the final
abilities, and the greedy choice also reads the held-out accuracies it is scored
against, so every value here is an upper bound rather than a rule a practitioner
can run. It marks how far selection could go, not how to get there.}

\aiadd{The pool is capped at 300 models drawn once with a fixed seed from the
cleaned training pool ($992$ models for IFEval~\cite{zhou2023ifeval}, $811$ for
MATH~\cite{hendrycks2021math}), and this capped pool is the reference: both the
full-pool correlation and the fitted abilities come from it. The held-out split
is the same fixed seed-7 $10\%$ test set used in the other
experiments~\cite{fourrier2024openllm}, $110$ IFEval models and $90$ MATH models,
so the curves line up with Sections~\ref{subsec:mcq-trainsize}
and~\ref{subsec:mcq-diversity}. We fit the abilities once, sort the 300 models
into five equal latent-ability groups of $60$, and report at $\mathrm{SE}\le0.3$.
On IFEval the groups span $\theta$ from about $-3.2$ to $1.6$
(Figure~\ref{fig:mcq-oracle-axis}). The full 300-model pool reaches $r=0.950$ on
IFEval and $r=0.933$ on MATH.}

\aiadd{The greedy oracle holds the reference correlation with very few models
(Table~\ref{tab:mcq-oracle}, Figure~\ref{fig:mcq-oracle}). On IFEval it reaches
$r=0.941$ at $N=50$, inside the tight $0.02$ band of the 300-model reference, and
$r=0.953$ at $N=100$, just above the reference. A random draw needs about $156$
models to come within $0.02$ of the reference and about $88$ to come within
$0.05$, so the greedy oracle holds the $0.02$ band with about $3.1\times$ fewer
models than random and the $0.05$ band with about $1.8\times$ fewer. We built the
greedy search on IFEval only. On MATH each calibration runs the long adaptive test
to $\mathrm{SE}\le0.3$ (about $164$ items), which makes the sequential greedy build
too costly for the compute budget, and the MATH fit is unstable enough (see below)
that a search reading held-out $r$ would chase noise.}

\aiadd{Spreading models across the ability groups is weaker and depends on the
benchmark. On IFEval the stratified oracle only matches random: it first comes
within $0.02$ of the reference at $N=300$, the whole pool, and within $0.05$ at
$N=150$, both worse than random at the same tolerance ($0.5$ and $0.6\times$). On
MATH it looks stronger, holding within $0.02$ at $N=100$, about $2.8\times$ fewer
than random (Figures~\ref{fig:mcq-oracle-math} and~\ref{fig:mcq-oracle-math-axis}),
but the MATH curve is noisy and its full-pool fit is unstable, with
the same 300 models re-fitting to $r$ between $0.90$ and $0.93$, so that factor is
not reliable. The stable reading is IFEval's: quantile spanning on its own is not
enough, and the dependable saving needs the greedy search that reads the held-out
correlation, which is exactly the information a deployable rule lacks.}

\aiadd{Read against the a-priori strategies in
Section~\ref{subsec:mcq-diversity}, the oracle fixes the two ends of the range.
Random sits at the bottom, needing about $156$ IFEval models for the $0.02$ band.
The greedy oracle sits at the top, holding the same band at $50$. The deployable
diversity strategies, which saved about $1.4$ to $1.6\times$, fall between the
two, closer to random than to the oracle. The practical reading is that most of
the theoretical saving depends on post-hoc ability information, so a real
selection rule should aim for the diversity end and treat the roughly $3\times$
figure as a ceiling, not a target.}

\aiadd{Several caveats bound these numbers. They are upper bounds by construction,
since the abilities come from the final calibration and the greedy search reads
the held-out correlation, so no a-priori rule can match them. The holding points
sit at the small-$N$ end, where a 3PL fit on $50$ to $100$ respondents against
several hundred items is noisy, so the min-$N$ values should be read as a band
rather than exact counts, and the savings factors come from interpolating the
random mean curve to the tolerance line. The greedy search on IFEval is also
restricted to the smallest grid sizes for tractability, built to $N=100$. MATH is
the least stable benchmark here: its adaptive test runs long, about $164$ items to
reach $\mathrm{SE}\le0.3$ against about $20$ on IFEval, its full-pool correlation
moves by about $0.03$ across re-fits of the same models, and we did not build its
greedy oracle, so the MATH results are random and stratified only and should be
read with caution.}

\begin{table}[t]
  \centering
  \caption{\aiadd{Minimum calibration models to hold the full-pool correlation at
  $\mathrm{SE}\le0.3$. ``Min $N$'' is the smallest pool size whose mean held-out
  $r$ stays within the tolerance of the full-pool reference; the random row gives
  the interpolated size where the random mean curve reaches the same line, and
  ``Save'' is that random size divided by the oracle Min $N$. Oracle rows use
  post-hoc fitted abilities (greedy also reads held-out $r$), so they are upper
  bounds. The greedy oracle was built on IFEval only; MATH is noisier and its
  full-pool $r$ moves across re-fits, so its stratified saving is not reliable.
  Source: \texttt{oracle\_subset\_selection/README.md} and
  \texttt{results.csv}.}}
  \label{tab:mcq-oracle}
  \small
  \begin{tabular}{l l r r r r}
    \toprule
     & & \multicolumn{2}{c}{within $0.02$} & \multicolumn{2}{c}{within $0.05$} \\
    \cmidrule(lr){3-4}\cmidrule(lr){5-6}
    Benchmark & Strategy & Min $N$ ($r$) & Save & Min $N$ ($r$) & Save \\
    \midrule
    IFEval & random ($r=0.950$)   & 156         & n/a         & 88         & n/a         \\
    IFEval & oracle-stratified    & 300 (0.950)& 0.5$\times$  & 150 (0.926)& 0.6$\times$  \\
    IFEval & oracle-greedy        & 50 (0.941) & 3.1$\times$  & 50 (0.941) & 1.8$\times$  \\
    \midrule
    MATH & random ($r=0.933$)     & 280         & n/a         & 199        & n/a         \\
    MATH & oracle-stratified      & 100 (0.921) & 2.8$\times$ & 100 (0.921)& 2.0$\times$  \\
    MATH & oracle-greedy          & \multicolumn{4}{c}{not run (see text)} \\
    \bottomrule
  \end{tabular}
\end{table}

\begin{figure}[t]
  \centering
  \includegraphics[width=\columnwidth]{figures/oracle_curve_ifeval.png}
  \caption{\aiadd{Held-out correlation against pool size on IFEval at
  $\mathrm{SE}\le0.3$. The shaded band is the random seed spread, the black line
  is the 300-model full-pool reference, and the circled markers show where each
  oracle strategy first holds within $0.02$ of it. The greedy oracle (red) holds
  the reference at $N=50$; the stratified oracle (blue) tracks random. Source:
  \texttt{figures/oracle\_curve\_ifeval.png}.}}
  \label{fig:mcq-oracle}
\end{figure}

\begin{figure}[t]
  \centering
  \includegraphics[width=\columnwidth]{figures/latent_axis_ifeval.png}
  \caption{\aiadd{The 300 IFEval calibration models placed along the oracle latent
  ability axis (full-pool EAP $\theta$), colored by the five equal-size groups.
  Triangles mark the theoretical-best subset (greedy, $N=50$), which spreads
  across the whole ability range. Source:
  \texttt{figures/latent\_axis\_ifeval.png}.}}
  \label{fig:mcq-oracle-axis}
\end{figure}

\begin{figure}[t]
  \centering
  \includegraphics[width=\columnwidth]{figures/oracle_curve_math.png}
  \caption{\aiadd{Held-out correlation against pool size on MATH at
  $\mathrm{SE}\le0.3$, with the random band, the full-pool reference line, and the
  oracle strategies. The greedy oracle is built only to $N=50$ here. Source:
  \texttt{figures/oracle\_curve\_math.png}.}}
  \label{fig:mcq-oracle-math}
\end{figure}

\begin{figure}[t]
  \centering
  \includegraphics[width=\columnwidth]{figures/latent_axis_math.png}
  \caption{\aiadd{The 300 MATH calibration models along the oracle latent ability
  axis (full-pool EAP $\theta$), colored by the five latent groups. No subset is
  marked because the greedy oracle was not built on MATH. Source:
  \texttt{figures/latent\_axis\_math.png}.}}
  \label{fig:mcq-oracle-math-axis}
\end{figure}

\aibullets{%
\begin{itemize}
  \item An oracle greedy subset holds the 300-model reference correlation within
  $0.02$ using $50$ IFEval models ($r=0.941$ against a full-pool $r=0.950$), about
  $3.1\times$ fewer than the roughly $156$ models random needs.
  \item On IFEval, spreading models across latent-ability groups (stratified) does
  not beat random; the dependable saving comes only from the greedy search that
  reads held-out $r$, so it is an upper bound, not a deployable rule.
  \item This oracle bounds the a-priori diversity result
  (Section~\ref{subsec:mcq-diversity}): random is the floor, the oracle is the
  ceiling, and practical diversity selection ($1.4$ to $1.6\times$) sits between
  them, closer to random.
  \item Caveats: upper bound by construction, noisy small-$N$ 3PL fits, greedy
  restricted to small $N$ for compute, and MATH runs a long full-pool test
  (about $164$ items), so its oracle points are the least stable.
\end{itemize}}
